\begin{cnabstract}
摘要应能独立成篇，让读者不必通读全文，也能掌握论文的核心信息。
其内容应覆盖与论文相当的主要信息，便于读者判断是否值得进一步阅读，
也便于二次文献引用。摘要通常要写明研究目的、方法、结果和结论，
其中结果和结论是重点。

\textcolor{red}{
说明：要求200$\sim$300字，“关键词”为小四号黑体，关键词的个数为3$\sim$8个。
关键词的排序，通常应按研究的对象、性质（问题）和采取的手段排序，而不应任意排列。
“关键词”后面不加冒号，空连个字符，关键词与关键词之间用全角逗号（，）隔开，
最后一个关键词后不加任何标点。
}

\end{cnabstract}

\begin{enabstract}
This is the English abstract of the thesis. It should summarise the research
problem, the method and the main conclusions of the paper.
\par
New paragraph of the English abstract.

\vspace{2em}
说明：\par 
Keywords 为小四号 Times New Roman，加粗。外文关键词（keyword）一律小写，
并与中文关键词一一对应。\textbf{Keywords} 后面不加冒号，空两个字符，
keyword 与 keyword 之间用半角逗号（,）隔开，每个（,）后空 1 个字符，
最后一个 keyword 后不加任何标点符号。
\par
\textcolor{red}{注意：此处行距已经设置为英文段落的间距
（即基础行距 1.2，行距设置为 1.5 倍行距）。}
\end{enabstract}
